\NormalFont\ShortTitle{1 Tessalonicenses}
{\MT Primeira Carta de Paulo aos Tessalonicenses

\par }\ChapOne{1}{\PP \VerseOne{1}Paulo, Silvano, e Timóteo, à igreja dos tessalonicenses, em Deus Pai, e no Senhor Jesus Cristo; haja convosco a graça e a paz de Deus, nosso Pai, e do Senhor Jesus Cristo.
\VS{2}Sempre damos graças a Deus por todos vós, fazendo menção de vós em nossas orações.
\VS{3}E, diante do nosso Deus e Pai, nós nos lembramos sem cessar da obra da vossa fé, do trabalho do amor, e da firmeza da esperança no nosso Senhor Jesus Cristo.
\VS{4}Sabemos, irmãos amados por Deus, que
{\ADD{por ele}} fostes escolhidos,
\FTNT{*}{{\FR 1:4 }
Lit. sabendo, irmãos amados por Deus, da vossa eleição
}
\VS{5}pois o nosso evangelho não foi até vós somente com palavras, mas também com poder, com o Espírito Santo, e com plena certeza; assim como sabeis de que maneira estivemos entre vós para o vosso benefício.
\VS{6}E vos tornastes imitadores nossos e do Senhor.
{\ADD{Mesmo}} em meio a muita aflição, recebestes a palavra com a alegria do Espírito Santo.
\VS{7}Dessa maneira vos tornastes referências a todos os crentes na Macedônia e Acaia.
\VS{8}Pois de vós ressoou a palavra do Senhor, não somente na Macedônia e Acaia, mas a vossa fé em Deus também se espalhou em todo lugar, de tal maneira que não precisamos falar coisa alguma.
\VS{9}Pois eles mesmos anunciam como fomos recebidos por vós,
\FTNT{†}{{\FR 1:9 }
como fomos recebidos por vós
Lit. que recepção tivemos até vós
} e como vos convertestes, deixando os ídolos, para servir ao Deus vivo e verdadeiro;
\VS{10}e para esperar dos céus o seu Filho, a quem ele ressuscitou dos mortos: Jesus, que nos livra da ira futura.

\par }\Chap{2}{\PP \VerseOne{1}Pois vós mesmos, irmãos, sabeis que a nossa passagem por entre vós não foi inútil.
\VS{2}Porém, mesmo que antes, em Filipos, tenhamos sofrido e sido maltratados, como sabeis, tivemos no nosso Deus ousadia para vos falar o evangelho de Deus em meio a muita oposição.
\VS{3}Pois a nossa exortação não
{\ADD{foi}} com engano, nem com impureza, nem fraude;
\VS{4}mas, como fomos aprovados por Deus para que o evangelho nos fosse confiado, assim falamos; não para agradar as pessoas, mas sim a Deus, que prova os nossos corações.
\VS{5}Pois, como sabeis, nunca usamos palavras de lisonja, nem pretexto de ganância; Deus é testemunha.
\VS{6}Não buscamos a glória humana, nem de vós, nem de outros, ainda que tínhamos autoridade, como apóstolos de Cristo, para demandar de vós.
\VS{7}Porém fomos suaves entre vós, como uma mãe que dá carinho aos seus filhos.
\VS{8}Assim nós, tendo tanta afeição por vós, queríamos de boa vontade compartilhar convosco não somente o evangelho de Deus, mas também as nossas próprias almas, porque éreis queridos por nós.
\VS{9}Pois vos lembrais, irmãos, do nosso trabalho e fadiga; porque enquanto vos pregávamos o evangelho de Deus, trabalhávamos noite e dia para que não fôssemos um peso para vós.
\VS{10}Vós e Deus
{\ADD{sois}} testemunhas de como foi santa, justa, e irrepreensível a maneira da qual nos comportamos convosco, que credes.
\VS{11}Assim como sabeis que, como um pai aos seus filhos, exortávamos, consolávamos, e testemunhávamos a cada um de vós,
\VS{12}para que andásseis de maneira digna diante de Deus, que vos chama para o seu reino e glória.
\VS{13}Por isso também agradecemos sem cessar a Deus, que, quando recebestes a palavra da Deus pregada por nós, a recebestes, não
{\ADD{como}} palavra de homens, mas (conforme em verdade é)
{\ADD{como}} a palavra de Deus, a qual também opera em vós, que credes.
\VS{14}Pois vós, irmãos, vos tornastes imitadores das igrejas de Deus que estão na Judeia em Cristo Jesus; porque também sofrestes as mesmas coisas dos vossos próprios compatriotas, como também eles dos judeus;
\VS{15}que também mataram o Senhor Jesus e os seus próprios profetas, e nos perseguiram, e não agradam a Deus, e são contrários a todos:
\VS{16}eles nos impedem de falar aos gentios, para que sejam salvos. Assim eles estão se enchendo constantemente de pecados, mas, por fim, ira veio sobre eles até o fim.
\FTNT{*}{{\FR 2:16 }
Ou: finalmente
}
\VS{17}Porém, irmãos, quando nos separamos de vós por algum tempo, em presença, não no coração, buscamos mais, com muito desejo, ver o vosso rosto;
\VS{18}por isso quisemos, pelo menos eu, Paulo, vos visitar uma vez e outra; mas Satanás nos impediu.
\VS{19}Pois, qual é nossa esperança, alegria, ou coroa de orgulho, diante do nosso Senhor Jesus Cristo na sua vinda? Acaso não sois vós?
\VS{20}Porque vós sois o nosso orgulho e alegria.

\par }\Chap{3}{\PP \VerseOne{1}Por isso, como não podíamos mais suportar, decidimos ficar sozinhos em Atenas,
\VS{2}e enviamos Timóteo, nosso irmão e servidor de Deus, e cooperador nosso no Evangelho de Cristo, para vos fortalecer e vos encorajar quanto a vossa fé;
\VS{3}a fim de que ninguém se abale por essas aflições; pois vós mesmos sabeis que para isso fomos destinados.
\FTNT{*}{{\FR 3:3 }
Ou: designados
}
\VS{4}Pois quando ainda estávamos convosco vós dizíamos com antecedência que seríamos afligidos, como aconteceu, e vós sabeis.
\VS{5}Por isso, como não podia mais suportar, eu
{\ADD{o}} enviei para saber de vossa fé.
{\ADD{Eu temia}} que o tentador houvesse vos tentado, e o nosso trabalho houvesse sido inútil.
\VS{6}Mas agora, depois que Timóteo voltou a nós da vossa presença, e nos trouxe boas notícias da vossa fé e amor, e de como sempre tendes boa lembrança de nós, e desejais muito nos ver, como também nós a vós;
\VS{7}por isso, irmãos, ficamos consolados acerca de vós em toda a nossa aflição e necessidade, pela vossa fé.
\VS{8}Pois agora vivemos, se vós estais firmes no Senhor.
\VS{9}Pois como agradecemos a Deus por vós, por toda a alegria com que nos alegramos por vossa causa, diante do nosso Deus!
\VS{10}E de noite e de dia com extrema intensidade oramos para que possamos ver o vosso rosto, e supramos o que falta na vossa fé.
\VS{11}E que o mesmo Deus, nosso Pai, e o nosso Senhor Jesus Cristo, preparem o nosso caminho até vós.
\VS{12}E o Senhor vos faça crescer e exceder em amor uns pelos outros, e por todos, à semelhança de nós por vós;
\VS{13}a fim de que fortaleça os vossos corações, irrepreensíveis em santidade diante do nosso Deus e Pai, na vinda do nosso Senhor Jesus Cristo com todos os seus santos.

\par }\Chap{4}{\PP \VerseOne{1}Portanto, irmãos, no restante, vos rogamos e exortamos no Senhor Jesus, que, assim como recebestes de nós como deveis andar e agradar a Deus, que façais
{\ADD{assim}} cada vez mais.
\VS{2}Pois vós sabeis quais mandamentos vos demos pelo Senhor Jesus.
\VS{3}Pois esta é a vontade de Deus: a vossa santificação, que não pratiqueis pecado sexual;
\VS{4}que cada um de vós saiba ser ter o seu instrumento
\FTNT{*}{{\FR 4:4 }
Não há um consenso sobre o verdadeiro sentido da palavra “instrumento” neste versículo. A primeira interpretação é a de que se refere ao corpo, isto é, “saiba ter domínio do corpo em santidade e honra”. Outra interpretação é a de que instrumento se refira à esposa, o que se traduziria: “saiba se casar com a sua esposa em santidade e honra”, porém essa interpretação é menos provável que a primeira.
} em santidade e honra;
\VS{5}não na paixão do desejo malicioso, como os gentios que não conhecem Deus.
\VS{6}Ninguém oprima ou engane ao seu irmão nisso, porque o Senhor é vingador de todas essas coisas, como também antes vos dissemos e demos testemunho.
\VS{7}Pois Deus não nos chamou para a impureza, mas sim para a santificação.
\VS{8}Portanto quem rejeita
{\ADD{isso}} , rejeita não ao ser humano, mas sim a Deus, que nos deu também o seu Espírito Santo.
\VS{9}Mas quanto ao amor fraternal, não precisais que eu vos escreva, pois vós mesmos estais instruídos por Deus que vos ameis uns aos outros.
\VS{10}Pois, de fato,
{\ADD{já}} fazeis assim a todos os irmãos de toda a Macedônia. Porém, irmãos, vos exortamos que façais{\ADD{isso}} ainda mais.
\VS{11}E procurai viver quietos, trantando dos vossos próprios assuntos, e trabalhando com as vossas próprias mãos, como já vos mandamos;
\VS{12}para que andeis de maneira respeitosa com os que estão de fora, e não necessiteis de nada.
\VS{13}Mas irmãos, não quero que desconheçais acerca dos que morreram,
\FTNT{†}{{\FR 4:13 }
Lit. acerca dos que dormem. também vv. 14, 15
} para que não vos entristeçais como os outros que não têm esperança.
\VS{14}Pois, se cremos que Jesus morreu e ressuscitou, assim também aos que morreram em Jesus, Deus os trará de volta com ele.
\VS{15}Pois dizemos isto pela palavra do Senhor, que nós, os que ficarmos vivos até a vinda do Senhor, não iremos antes dos que morreram.
\VS{16}Pois o mesmo Senhor descerá do céu com grande aclamação, com voz de arcanjo, e com trombeta de Deus; e os que morreram em Cristo ressuscitarão primeiro;
\VS{17}Em seguida nós, os que ficarmos vivos, seremos arrebatados juntamente com eles nas nuvens, para encontrar o Senhor no ar, e assim estaremos com o Senhor para sempre.
\VS{18}Portanto consolai-vos uns aos outros com essas palavras.

\par }\Chap{5}{\PP \VerseOne{1}Mas irmãos, acerca dos tempos e das estações, não necessitais de que eu vos escreva;
\VS{2}pois vós mesmos bem sabeis que o dia do Senhor virá como um ladrão de noite.
\VS{3}Pois quando disserem: “Paz e segurança”, então virá sobre eles repentina destruição, como as dores de parto da grávida, e de maneira nenhuma escaparão.
\VS{4}Mas vós, irmãos, não estais em trevas, para que aquele dia vos tome de surpresa como um ladrão.
\VS{5}Todos vós sois filhos da luz e filhos do dia; não somos da noite nem das trevas.
\VS{6}Portanto, não durmamos, como os outros; em vez disso, vigiemos e sejamos sóbrios.
\VS{7}Pois os que dormem, dormem de noite; e os que ficam bêbados embebedam-se de noite.
\VS{8}Mas nós, que somos do dia, sejamos sóbrios, e nos vistamos da couraça da fé e do amor, e do capacete da esperança da salvação.
\VS{9}Porque Deus não nos destinou para a ira, mas, sim, para obtermos a salvação, por meio do nosso Senhor Jesus Cristo,
\VS{10}que morreu por nós, a fim de que, se estivermos vigiando ou dormindo, nós vivamos juntamente com ele.
\VS{11}Portanto exortai-vos uns aos outros, e edificai-vos uns aos outros, como já fazeis.
\VS{12}Irmãos, nós vos rogamos que reconheçais os que trabalham entre vós, e que lideram sobre vós no Senhor, e vos advertem;
\VS{13}e os estimai muito com amor, por causa do trabalho deles. Tende paz entre vós.
\VS{14}Rogamo-vos também, irmãos, que advirtais os desordeiros, consoleis os de pouco ânimo, apoieis aos fracos, e tenhais paciência com todos.
\VS{15}Tende cuidado para que ninguém retribua o mal com o mal, mas sempre segui o bem, tanto uns com os outros, como com todos.
\VS{16}Alegrai-vos sempre.
\VS{17}Orai sem cessar.
\VS{18}Agradecei em tudo, pois essa é a vontade de Deus em Cristo Jesus para vós mesmos.
\VS{19}Não apagueis o Espírito.
\VS{20}Não desprezeis as profecias.
\VS{21}Examinai tudo, e mantende o que é bom.
\VS{22}Não haja entre vós qualquer forma de mal.
\VS{23}E o próprio Deus da paz vos santifique em tudo; e que todo o vosso espírito, alma e corpo sejam preservados irrepreensíveis na vinda do nosso Senhor Jesus Cristo.
\VS{24}Fiel é aquele que vos chama, e ele também
{\ADD{o}} fará.
\VS{25}Irmãos, orai por nós.
\VS{26}Saudai a todos os irmãos com beijo santo.
\VS{27}Eu vos mando pelo Senhor que esta carta seja lida a todos os santos irmãos.
\VS{28}A graça de nosso Senhor Jesus Cristo
{\ADD{seja}} convosco. Amém.
{\ADD{A primeira carta aos tessalonicenses foi escrita de Atenas}}
\par }